\chapter{General Relativity: Setup and Conventions}
\label{app:GR_basics}

This appendix collects the geometric conventions used in the project.
The purpose is not to derive the Einstein equations, but to fix notation
and sign conventions so that the definitions used in the main text are
unambiguous.

We use natural units with
\[
c=G=1,
\]
and the metric signature is chosen to be
\[
(+,-,-,-).
\]
With this convention, timelike four-velocities satisfy
\[
u^\mu u_\mu = 1 .
\]
Greek indices \(\mu,\nu,\ldots\) run over spacetime coordinates
\(0,1,2,3\), while Latin indices \(i,j,\ldots\) denote spatial
components \(1,2,3\). Symmetrization and antisymmetrization are defined
with unit weight,
\[
A_{(\mu\nu)}
=
\frac{1}{2}
\left(A_{\mu\nu}+A_{\nu\mu}\right),
\qquad
A_{[\mu\nu]}
=
\frac{1}{2}
\left(A_{\mu\nu}-A_{\nu\mu}\right).
\]

The metric tensor \(g_{\mu\nu}\) defines the line element
\begin{equation}
ds^2 = g_{\mu\nu}\,dx^\mu dx^\nu .
\end{equation}
Indices are raised and lowered with \(g^{\mu\nu}\) and \(g_{\mu\nu}\),
where
\[
g^{\mu\alpha}g_{\alpha\nu}
=
\delta^\mu{}_\nu .
\]
The determinant of the metric is denoted
\[
g \equiv \det(g_{\mu\nu}) .
\]
For the mostly-minus signature used here, \(g<0\) for ordinary
Lorentzian coordinate systems, and the invariant spacetime volume
element is
\[
\sqrt{-g}\,d^4x .
\]

All covariant derivatives are defined using the Levi--Civita connection,
which is torsion free and metric compatible. Acting on contravariant and
covariant vectors, the covariant derivative is
\begin{equation}
\nabla_\mu V^\nu
=
\partial_\mu V^\nu
+
\Gamma^\nu{}_{\mu\lambda}V^\lambda,
\qquad
\nabla_\mu V_\nu
=
\partial_\mu V_\nu
-
\Gamma^\lambda{}_{\mu\nu}V_\lambda .
\end{equation}
The Christoffel symbols are
\begin{equation}
\Gamma^\rho{}_{\mu\nu}
=
\frac{1}{2}g^{\rho\sigma}
\left(
\partial_\mu g_{\nu\sigma}
+
\partial_\nu g_{\mu\sigma}
-
\partial_\sigma g_{\mu\nu}
\right).
\label{eq:christoffel_def}
\end{equation}

The curvature convention is fixed by defining the Riemann tensor as
\begin{equation}
R^\rho{}_{\sigma\mu\nu}
=
\partial_\mu \Gamma^\rho{}_{\nu\sigma}
-
\partial_\nu \Gamma^\rho{}_{\mu\sigma}
+
\Gamma^\rho{}_{\mu\lambda}\Gamma^\lambda{}_{\nu\sigma}
-
\Gamma^\rho{}_{\nu\lambda}\Gamma^\lambda{}_{\mu\sigma}.
\label{eq:riemann_def}
\end{equation}
The Ricci tensor and Ricci scalar are obtained by contraction,
\begin{equation}
R_{\mu\nu}
=
R^\rho{}_{\mu\rho\nu},
\qquad
R
=
g^{\mu\nu}R_{\mu\nu}.
\label{eq:ricci_def}
\end{equation}
With these definitions, the Einstein tensor is
\begin{equation}
G_{\mu\nu}
=
R_{\mu\nu}
-
\frac{1}{2}R\,g_{\mu\nu}.
\label{eq:einstein_tensor_def}
\end{equation}
For the Levi--Civita connection it satisfies the twice-contracted
Bianchi identity,
\begin{equation}
\nabla^\mu G_{\mu\nu}=0 .
\end{equation}

Matter enters through the energy--momentum tensor \(T_{\mu\nu}\), which
encodes energy density, momentum density, and stresses. Local
energy--momentum conservation is written as
\begin{equation}
\nabla^\mu T_{\mu\nu}=0 .
\end{equation}
For a perfect fluid in the signature convention used here,
\begin{equation}
T_{\mu\nu}
=
(\epsilon+p)u_\mu u_\nu
-
p\,g_{\mu\nu},
\label{eq:perfect_fluid_Tmn}
\end{equation}
where \(\epsilon\) is the energy density, \(p\) is the pressure, and
\(u^\mu\) is the fluid four-velocity.

The Einstein field equations are written as
\begin{equation}
R_{\mu\nu}
-
\frac{1}{2}R\,g_{\mu\nu}
=
8\pi T_{\mu\nu}.
\label{eq:einstein_equation}
\end{equation}
Since \(c=G=1\), the gravitational coupling is \(8\pi\). The sign
conventions in this appendix are consistent with the mostly-minus metric
signature. When comparing with sources that use the mostly-plus
signature \((- ,+,+,+)\), several intermediate signs may differ, although
coordinate-invariant physical statements are unchanged.
